% ============================================================
%  Pre-Analysis Feasibility Assessment in Bayesian Inference
%  A practical guide for study design and data requirements
% ============================================================
\documentclass[11pt, a4paper]{article}

\usepackage[margin=2.5cm]{geometry}
\usepackage{amsmath, amssymb, amsthm}
\usepackage{booktabs}
\usepackage{array}
\usepackage{xcolor}
\usepackage{hyperref}
\usepackage{parskip}
\usepackage{enumitem}
\usepackage{mdframed}
\usepackage{titlesec}
\usepackage{microtype}

\hypersetup{
    colorlinks = true,
    linkcolor  = blue!60!black,
    urlcolor   = blue!60!black,
    citecolor  = blue!60!black,
}

\definecolor{boxbg}{RGB}{245, 248, 255}
\definecolor{boxborder}{RGB}{70, 130, 180}

% Highlighted box for key results
\newmdenv[
    backgroundcolor = boxbg,
    linecolor       = boxborder,
    linewidth       = 1.2pt,
    innerleftmargin = 12pt,
    innerrightmargin= 12pt,
    innertopmargin  = 10pt,
    innerbottommargin=10pt,
    skipabove       = 10pt,
    skipbelow       = 10pt,
]{keybox}

\titleformat{\section}{\large\bfseries}{{\thesection.}}{0.6em}{}[\vspace{-2pt}\rule{\linewidth}{0.4pt}]
\titleformat{\subsection}{\normalsize\bfseries}{{\thesubsection}}{0.6em}{}

% ── document ────────────────────────────────────────────────────────────────
\begin{document}

\begin{center}
    {\LARGE \textbf{Pre-Analysis Feasibility Assessment}}\\[6pt]
    {\large in Bayesian Inference and Time Series Modelling}\\[10pt]
    {\normalsize A practical guide for deciding whether a study is worth running}\\[4pt]
    {\small \textit{Notes compiled from analysis of AR(1) variability inference in AGN X-ray light curves}}
\end{center}

\vspace{4pt}
\noindent\rule{\linewidth}{0.8pt}
\vspace{6pt}

\tableofcontents
\vspace{10pt}
\noindent\rule{\linewidth}{0.4pt}

% ──────────────────────────────────────────────────────────────────────────────
\section{The Core Principle: The Information Budget}

Every inference problem, regardless of method, is fundamentally a question of whether the
planned dataset contains enough \textit{information} to constrain the parameters of interest
to the required precision.

\begin{keybox}
\textbf{Central question (ask this before collecting a single data point):}

\medskip
\textit{``Does my planned dataset have sufficient information content to answer my scientific
question to the required precision?''}

\medskip
If the answer is no, no amount of computational sophistication --- Hamiltonian Monte Carlo,
variational inference, or otherwise --- can compensate for missing information.
\end{keybox}

The information available in a dataset is set by three factors that must \emph{all} be
satisfied simultaneously:
\begin{enumerate}[noitemsep]
    \item \textbf{Quantity:} how many effectively independent measurements exist.
    \item \textbf{Quality:} how precisely each measurement constrains the parameter of interest.
    \item \textbf{Identifiability:} whether the parameter is uniquely determined by the data
          at all, or whether it is degenerate with other parameters.
\end{enumerate}


% ──────────────────────────────────────────────────────────────────────────────
\section{Time Series Analysis: the AR(1) Case}

\subsection{Effective sample size is the only sample size that matters}

For any correlated time series, the raw length $T$ is misleading as a measure of
information content. Adjacent points are correlated, so they carry redundant information.
For an AR(1) process $X_t = \phi\, X_{t-1} + \varepsilon_t$ with autocorrelation
parameter $\phi = e^{-1/\tau}$, the effective number of statistically independent
observations is:

\begin{equation}
    \boxed{
        n_\mathrm{eff} = T \cdot \frac{1 - \phi}{1 + \phi}
                       = T \cdot \frac{1 - e^{-1/\tau}}{1 + e^{-1/\tau}}
    }
    \label{eq:neff}
\end{equation}

This quantity is \textbf{fully computable before any analysis}, requiring only a prior
belief (or order-of-magnitude estimate) of $\tau$.

\medskip
Table~\ref{tab:neff_examples} shows $n_\mathrm{eff}$ for three representative
timescales with $T = 100$ steps, illustrating why inference quality degrades sharply
as $\tau$ approaches or exceeds $T$.

\begin{table}[h!]
\centering
\caption{Effective sample size for $T=100$ as a function of true timescale $\tau$.}
\label{tab:neff_examples}
\begin{tabular}{ccccc}
\toprule
True $\tau$ & $\phi = e^{-1/\tau}$ & $n_\mathrm{eff}$ & Correlation blocks & Inference quality \\
\midrule
$2$  & $0.607$ & $24.5$ & $\sim 50$ & Well constrained \\
$38$ & $0.974$ & $1.3$  & $\sim 2.6$ & Barely constrained \\
$50$ & $0.980$ & $1.0$  & $\sim 2$   & Essentially unconstrained \\
\bottomrule
\end{tabular}
\end{table}

\subsection{Minimum time series length as a function of precision target}

Using equation~\eqref{eq:neff} and the empirical rule that reliable estimation of
a parameter requires $n_\mathrm{eff} \gtrsim 10$--$20$, we can invert to obtain the
required time series length:

\begin{equation}
    T_\mathrm{min} \approx n_\mathrm{eff}^\mathrm{target} \cdot \frac{1 + e^{-1/\tau}}{1 - e^{-1/\tau}}
    \approx n_\mathrm{eff}^\mathrm{target} \cdot 2\tau \quad (\text{for large } \tau)
\end{equation}

The practical rule of thumb expressed as a precision target on $\tau$ itself is:

\begin{table}[h!]
\centering
\caption{Required time series length as a function of precision target on $\tau$.}
\label{tab:Tmin}
\begin{tabular}{lc}
\toprule
Precision goal on $\hat{\tau}$ & Required time series length \\
\midrule
Detect that $\tau$ is finite (exclude $\tau \to \infty$) & $T \gtrsim 3\tau$ \\
Constrain $\tau$ to $\pm 50\%$ & $T \gtrsim 10\tau$ \\
Constrain $\tau$ to $\pm 20\%$ & $T \gtrsim 50\tau$ \\
\bottomrule
\end{tabular}
\end{table}

\begin{keybox}
\textbf{Design implication:} Before observing, estimate the largest timescale $\tau_\mathrm{max}$
you wish to measure. Ensure $T \gg \tau_\mathrm{max}$. If telescope or observing time
constraints prevent this, accept that $\tau$ cannot be tightly constrained and adjust
the scientific claim accordingly.
\end{keybox}

\subsection{Why more photons do not help when $T \lesssim \tau$}

A crucial distinction arises between two types of uncertainty:

\begin{enumerate}
    \item \textbf{Photon noise (quality):} Insufficient photons per time step means
    the spectrum cannot be fitted precisely. \emph{More photons directly help here.}

    \item \textbf{Temporal structure (quantity):} When $T \lesssim \tau$, the
    autocorrelation function has barely decayed over the full length of the time series.
    Even with infinite photons per step, you are observing only one or two realizations
    of the correlated process. \emph{More photons cannot substitute for more time steps.}
\end{enumerate}

Concretely, for $\tau = 50$ and $T = 100$, the full time series spans only two
correlation lengths. A single unlucky realisation --- where the process happens to stay
elevated for 70 steps then drop --- is statistically indistinguishable from $\tau = 200$
or even $\tau \to \infty$. The posterior spans the entire prior range regardless of
photon count.

\subsection{The signal-to-noise floor and minimum detectable amplitude}

Even if $n_\mathrm{eff}$ is adequate, the variability amplitude $\sigma$ of the
process must exceed the measurement noise. The minimum detectable amplitude is
approximately:

\begin{equation}
    \sigma_\mathrm{min} \sim \frac{\text{noise per step}}{\sqrt{n_\mathrm{eff}}}
\end{equation}

Both the quantity and quality constraints must be satisfied simultaneously.

% -----------------------------------------------------------------------
\subsection{Deriving $n_\mathrm{eff}$ from first principles}
\label{subsec:neff_derivation}

The formula in equation~\eqref{eq:neff} is not an empirical rule of thumb.
It follows from a classical result in time series statistics.
This section derives it carefully, starting from the only distributional assumption
required.

\subsubsection*{What distributional assumption is needed?}

No assumption of Gaussianity is required.
The derivation uses only \textbf{weak (second-order) stationarity}:
\begin{align}
    \mathrm{E}[X_t] &= \mu \quad \text{(constant mean)}, \\
    \mathrm{Var}(X_t) &= \sigma^2 \quad \text{(constant marginal variance)}, \\
    \mathrm{Cov}(X_t, X_{t+k}) &= \gamma(k) = \sigma^2\rho(k)
        \quad \text{(covariance depends only on lag } k\text{)}.
\end{align}
The formula for $n_\mathrm{eff}$ holds for any weakly stationary process.
Gaussianity is only required if you wish to make probability statements about the
distribution of the sample mean $\bar{X}$ itself; the variance formula is
purely a second-order (covariance) result.

\subsubsection*{Why $\mathrm{Var}(\bar{X})$ and not $\mathrm{Var}(X)$?}

$\mathrm{Var}(X_t) = \sigma^2$ is a fixed property of the process:
the marginal spread of a single observation.
It does not change with $T$ and carries no information about how precisely
anything can be estimated from a sample of size $T$.

$\mathrm{Var}(\bar{X})$ measures the \textbf{precision of the sample mean as an
estimator of $\mu$}.
This is what shrinks as $T$ grows, and the rate of shrinkage is the object that
reveals how much independent information the correlated observations provide.
More precisely, $n_\mathrm{eff}$ answers the question:
\begin{quote}
    \textit{How many i.i.d.\ draws would give the same estimation precision
    as $T$ correlated draws?}
\end{quote}
Precision is measured by $\mathrm{Var}(\bar{X})$, not by $\mathrm{Var}(X)$.
Hence $\mathrm{Var}(\bar{X})$ is the central object.

\subsubsection*{Exact derivation}

Define the sample mean:
\begin{equation}
    \bar{X} = \frac{1}{T}\sum_{t=1}^{T} X_t.
\end{equation}
Its variance expands as:
\begin{equation}
    \mathrm{Var}(\bar{X})
    = \frac{1}{T^2}\,\mathrm{Var}\!\left(\sum_{t=1}^T X_t\right)
    = \frac{1}{T^2}\sum_{t=1}^T\sum_{s=1}^T \mathrm{Cov}(X_t,\,X_s)
    = \frac{1}{T^2}\sum_{t=1}^T\sum_{s=1}^T \gamma(|t-s|).
\end{equation}
Now count how many pairs $(t,s)$ share each lag value $k = |t-s|$:
at lag $k=0$ there are $T$ pairs; at lag $k \geq 1$ there are $2(T-k)$ pairs.
Therefore:
\begin{equation}
    \sum_{t=1}^T\sum_{s=1}^T \gamma(|t-s|)
    = T\,\gamma(0) + 2\sum_{k=1}^{T-1}(T-k)\,\gamma(k).
\end{equation}
Dividing by $T^2$ and substituting $\gamma(0) = \sigma^2$,
$\gamma(k) = \sigma^2\rho(k)$:
\begin{equation}
    \boxed{
        \mathrm{Var}(\bar{X})
        = \frac{\sigma^2}{T}\!\left[1 + 2\sum_{k=1}^{T-1}\!\left(1-\frac{k}{T}\right)\rho(k)\right].
    }
    \label{eq:var_xbar_exact}
\end{equation}
The factor $(1 - k/T)$ is the \textbf{Bartlett triangular taper} (Bartlett 1946):
it arises because at lag $k$ only $T-k$ pairs are available, not $T$.

\subsubsection*{Large-$T$ simplification and the integrated autocorrelation time}

For large $T$, the term $k/T \to 0$ for all lags that contribute appreciably to the
sum, so equation~\eqref{eq:var_xbar_exact} simplifies to:
\begin{equation}
    \mathrm{Var}(\bar{X}) \approx \frac{\sigma^2}{T}\!\left[1 + 2\sum_{k=1}^{\infty}\rho(k)\right]
    = \frac{2\,\sigma^2\,\tau_\mathrm{int}}{T},
\end{equation}
where the \textbf{integrated autocorrelation time} is defined as:
\begin{equation}
    \tau_\mathrm{int} = \frac{1}{2} + \sum_{k=1}^{\infty}\rho(k).
    \label{eq:tauint}
\end{equation}
The variance formula~\eqref{eq:var_xbar_exact} originates with Bartlett (1946).
The integrated autocorrelation time $\tau_\mathrm{int}$ and the
$n_\mathrm{eff} = T/(2\tau_\mathrm{int})$ terminology were established in the
MCMC literature by Geyer (1992) and are the basis of the convergence diagnostics
in Gelman et al.\ (2013).
The analogous concept for correlated samples in survey sampling was introduced
earlier by Kish (1965).

\subsubsection*{Defining $n_\mathrm{eff}$ by comparison to i.i.d.\ data}

For i.i.d.\ data, $\rho(k) = 0$ for $k \geq 1$, so
$\mathrm{Var}(\bar{X})_\mathrm{iid} = \sigma^2 / n_\mathrm{eff}$.
Setting this equal to the correlated result:
\begin{equation}
    \frac{\sigma^2}{n_\mathrm{eff}} = \frac{2\,\sigma^2\,\tau_\mathrm{int}}{T}
    \quad\Longrightarrow\quad
    \boxed{n_\mathrm{eff} = \frac{T}{2\,\tau_\mathrm{int}}}.
    \label{eq:neff_tauint}
\end{equation}

\subsubsection*{Evaluating for AR(1)}

For the AR(1) process, $\rho(k) = \phi^{|k|} = e^{-|k|/\tau}$.
Summing the geometric series:
\begin{equation}
    \sum_{k=1}^{\infty} e^{-k/\tau}
    = \frac{e^{-1/\tau}}{1 - e^{-1/\tau}}.
\end{equation}
So $\tau_\mathrm{int} = \frac{1}{2} + \frac{e^{-1/\tau}}{1-e^{-1/\tau}}
= \frac{1}{1 - e^{-1/\tau}} - \frac{1}{2}$.
After simplification, $n_\mathrm{eff}$ becomes:
\begin{equation}
    n_\mathrm{eff} = T \cdot \frac{1 - e^{-1/\tau}}{1 + e^{-1/\tau}},
\end{equation}
which is exactly equation~\eqref{eq:neff}.
For large $\tau$, $e^{-1/\tau} \approx 1 - 1/\tau$, giving
$n_\mathrm{eff} \approx T/(2\tau)$,
and $\tau_\mathrm{int} \approx \tau$: the process timescale and the integrated
autocorrelation time are the same object.

% -----------------------------------------------------------------------
\subsection{Connection to Fisher information and posterior width}
\label{subsec:neff_fisher}

The result $n_\mathrm{eff} \approx T/(2\tau)$ is not merely a precision bound on
the sample mean.
It directly controls the posterior width for $\tau$ itself, through the Fisher
information.

\subsubsection*{The Cram\'er--Rao bound and Bernstein--von Mises theorem}

The Fisher information for parameter $\theta$ from a dataset of length $T$ is:
\begin{equation}
    \mathcal{I}_T(\theta)
    = -\,\mathrm{E}\!\left[\frac{\partial^2 \ln p(\mathbf{X}\mid\theta)}{\partial\theta^2}\right].
\end{equation}
The \textbf{Cram\'er--Rao bound} states that for any unbiased estimator $\hat\theta$:
\begin{equation}
    \mathrm{Var}(\hat\theta) \geq \frac{1}{\mathcal{I}_T(\theta)}.
\end{equation}
The \textbf{Bernstein--von Mises theorem} states that for large $T$,
regardless of the prior, the Bayesian posterior concentrates as:
\begin{equation}
    p(\theta \mid \mathbf{X}) \;\xrightarrow{T\to\infty}\;
    \mathcal{N}\!\left(\hat\theta_\mathrm{MLE},\;\mathcal{I}_T(\theta)^{-1}\right).
\end{equation}
The posterior variance equals the Cram\'er--Rao bound.
\textbf{Posterior width is therefore governed entirely by Fisher information.}

\subsubsection*{Fisher information for $\phi$ in a Gaussian AR(1)}

For a Gaussian AR(1), the conditional log-likelihood given $X_1$ is:
\begin{equation}
    \ln p(\mathbf{X}\mid\mu,\phi,\sigma^2)
    = -\frac{T-1}{2}\ln\sigma^2
    - \frac{1}{2\sigma^2}\sum_{t=2}^{T}\bigl[X_t - \mu - \phi(X_{t-1}-\mu)\bigr]^2.
\end{equation}
Taking the expected squared score with respect to $\phi$
(and using the fact that innovations $V_t$ are independent of $X_{t-1}$):
\begin{equation}
    \mathcal{I}_T(\phi)
    = \frac{1}{\sigma^2}\sum_{t=2}^{T}\mathrm{E}\bigl[(X_{t-1}-\mu)^2\bigr]
    = \frac{T-1}{\sigma^2}\cdot\frac{\sigma^2}{1-\phi^2}
    = \frac{T-1}{1-\phi^2}.
    \label{eq:fisher_phi}
\end{equation}
This scales as $T$ because each innovation $V_t$ is genuinely i.i.d.,
providing one independent piece of information about $\phi$ per time step.

\subsubsection*{Fisher information for $\tau$ via the chain rule}

Since $\phi = e^{-1/\tau}$, the chain rule gives:
\begin{equation}
    \mathcal{I}_T(\tau) = \mathcal{I}_T(\phi)\cdot\left(\frac{d\phi}{d\tau}\right)^{\!2},
    \qquad \frac{d\phi}{d\tau} = \frac{e^{-1/\tau}}{\tau^2} = \frac{\phi}{\tau^2}.
\end{equation}
Substituting equation~\eqref{eq:fisher_phi}:
\begin{equation}
    \mathcal{I}_T(\tau)
    = \frac{T-1}{1-\phi^2}\cdot\frac{\phi^2}{\tau^4}.
\end{equation}
In the large-$\tau$ limit ($\phi \to 1$),
use $1-\phi^2 = (1-\phi)(1+\phi) \approx \frac{2}{\tau}$ and $\phi^2 \approx 1$:
\begin{equation}
    \boxed{
        \mathcal{I}_T(\tau) \approx \frac{T}{2\tau^3}.
    }
    \label{eq:fisher_tau}
\end{equation}

\subsubsection*{Connecting $\mathcal{I}_T(\tau)$ to $n_\mathrm{eff}$}

The Cram\'er--Rao lower bound on the relative precision of $\hat\tau$ is:
\begin{equation}
    \frac{\mathrm{std}(\hat\tau)}{\tau}
    \geq \sqrt{\frac{1}{\tau^2\,\mathcal{I}_T(\tau)}}
    = \sqrt{\frac{2\tau}{T}}
    = \frac{1}{\sqrt{n_\mathrm{eff}}}.
    \label{eq:crb_tau_neff}
\end{equation}
The \textbf{relative posterior width for $\tau$ is lower-bounded by
$1/\sqrt{n_\mathrm{eff}}$}.
This is the formal statement connecting the informational concept ($n_\mathrm{eff}$),
the frequentist bound (Cram\'er--Rao), and the Bayesian posterior (Bernstein--von Mises)
into a single inequality.

For the case $\tau = 38$, $T = 100$: $n_\mathrm{eff} \approx 1.3$, so
$\mathrm{std}(\hat\tau)/\tau \geq 0.88$.
The standard deviation of the posterior for $\tau$ is nearly as large as $\tau$
itself---a genuinely diffuse posterior, not a numerical artefact.

\subsubsection*{Why $\mathcal{I}_T(\tau)$ is small: the reparameterisation cost}

The Fisher information for $\phi$ scales as $T$---fast accumulation.
But converting from $\phi$ (which lives near 1 for large $\tau$) to $\tau$ costs
a factor of $(d\phi/d\tau)^2 = \phi^2/\tau^4 \approx \tau^{-4}$.
Combining:
\begin{equation}
    \mathcal{I}_T(\tau)
    = \underbrace{\frac{T-1}{1-\phi^2}}_{\displaystyle\text{info about }\phi}
    \times
    \underbrace{\frac{\phi^2}{\tau^4}}_{\displaystyle\text{reparameterisation cost}}.
\end{equation}
A unit change in $\tau$ at $\tau = 38$ corresponds to a change in $\phi$ of only
$\phi/\tau^2 \approx 1/1444$.
The data can detect changes in $\phi$ readily, but the map from $\phi$-precision to
$\tau$-precision requires a large derivative---which is tiny when $\tau$ is large.
Table~\ref{tab:fisher_tau} summarises the resulting bounds numerically.

\begin{table}[h!]
\centering
\caption{Cram\'er--Rao lower bound on the relative standard deviation of $\hat\tau$
         for $T = 100$ at three timescales, compared to the simulation results.}
\label{tab:fisher_tau}
\begin{tabular}{ccccc}
\toprule
$\tau$ & $n_\mathrm{eff}$ & CRB: $\mathrm{std}(\hat\tau)/\tau \geq$ &
Posterior 5--95\% CI & CI width / $\tau$ \\
\midrule
$2$  & $24.5$ & $0.20$ & $[1.6,\;2.5]$    & $0.45$ \\
$38$ & $1.3$  & $0.88$ & $[7.1,\;77.5]$   & $1.85$ \\
$50$ & $1.0$  & $1.0$  & $[\lesssim 5,\;\gtrsim 60]$ & $>1$ \\
\bottomrule
\end{tabular}
\end{table}

The full chain from $n_\mathrm{eff}$ to posterior width is:
\begin{equation}
    n_\mathrm{eff} = \frac{T}{2\tau}
    \;\longrightarrow\;
    \mathcal{I}_T(\tau) = \frac{n_\mathrm{eff}}{\tau^2}
    \;\longrightarrow\;
    \text{posterior std} \approx \frac{1}{\sqrt{\mathcal{I}_T(\tau)}}
    = \frac{\tau}{\sqrt{n_\mathrm{eff}}}.
\end{equation}
When $n_\mathrm{eff} \ll 1$, $\mathcal{I}_T(\tau)$ is small, the CRB is large, and the
posterior is wide.
These are three descriptions of the same fact.

% -----------------------------------------------------------------------
\subsection{Statistical $n_\mathrm{eff}$ versus MCMC $n_\mathrm{eff}$}
\label{subsec:neff_two_kinds}

The term ``effective sample size'' appears in two completely different contexts in
Bayesian computation.
They measure different things, have different remedies, and must not be conflated.

\begin{table}[h!]
\centering
\caption{The two senses of effective sample size.}
\label{tab:neff_comparison}
\begin{tabular}{p{3.5cm} p{5.5cm} p{5.5cm}}
\toprule
 & \textbf{Statistical $n_\mathrm{eff}$} & \textbf{MCMC $n_\mathrm{eff}$} \\
\midrule
\textbf{Question answered} &
How much independent information does the observed data $\{C_t\}$ contain about
the parameter of interest? &
How many of the posterior draws produced by the sampler are effectively independent
of each other? \\[6pt]
\textbf{Determined by} &
The autocorrelation structure of the process: $\tau$ and $T$ &
The geometry of the posterior, step size, trajectory length, and sampler algorithm \\[6pt]
\textbf{Formula} &
$n_\mathrm{eff} = T/(2\tau_\mathrm{int})$ &
$n_\mathrm{eff}^\mathrm{MCMC} = N_\mathrm{draws}/\bigl(1 + 2\sum_{k}\hat\rho_k^\mathrm{chain}\bigr)$ \\[6pt]
\textbf{Fixed by} &
Collecting more data (longer light curve, $T \gg \tau$) &
Running more iterations, better step size, reparameterisation, better sampler \\[6pt]
\textbf{Role in the AR(1) problem} &
Sets the irreducible floor on how precisely $\tau$ can be estimated.
No sampler can overcome an insufficiently small statistical $n_\mathrm{eff}$ &
Reports how many independent posterior samples the chain has produced.
NumPyro's \texttt{summary()} outputs this (e.g.\ $n_\mathrm{eff}^\mathrm{MCMC} = 1338$
for $\hat\tau_{N_H}$ at $\tau=38$) \\
\bottomrule
\end{tabular}
\end{table}

\subsubsection*{How they are connected through the posterior geometry}

The AR(1) prior $p(\{X_t\} \mid \tau, \sigma)$ induces correlations between adjacent
latent states $X_1, X_2, \ldots, X_T$.
For large $\tau$, these correlations are strong: $X_1$ and $X_{50}$ still have
correlation $e^{-50/38} \approx 0.27$.
The 100 latent parameters $\{X_t\}$ do not span 100 independent directions in
parameter space; they span a low-dimensional ridge with intrinsic dimension
$\approx n_\mathrm{eff} \approx 1$--$2$.

HMC must navigate this ridge.
The leapfrog integrator's step size is constrained by the stiffest direction
(across the ridge), while progress along the ridge is slow.
This slows the \emph{MCMC} exploration of the posterior.
At the same time, the statistical $n_\mathrm{eff} \approx 1.3$ means the posterior
for $\tau$ is genuinely wide in the first place---even a perfect sampler with
MCMC $n_\mathrm{eff} = \infty$ would recover the same wide credible interval.

The same physical structure---long $\tau$ relative to $T$---therefore creates
problems for \emph{both} information content and sampler geometry, but through
entirely different mechanisms.

\subsubsection*{A concrete illustration for the AGN problem}

At $\tau = 38$, $T = 100$:
\begin{itemize}[noitemsep]
    \item Statistical $n_\mathrm{eff} \approx 1.3$: the data contain fewer than
          two independent correlation lengths.
          The Fisher information bound gives
          $\mathrm{std}(\hat\tau)/\tau \geq 0.88$: the posterior
          \emph{must} be wide, by physics.
    \item MCMC $n_\mathrm{eff}^\mathrm{MCMC} = 1338$ (from NumPyro \texttt{summary()}
          after 6000 draws): the sampler is mapping the posterior well, producing
          1338 effectively independent draws of the wide posterior.
\end{itemize}
Running ten times more HMC iterations would improve MCMC $n_\mathrm{eff}$ to
$\sim 13\,000$ but would \emph{not} narrow the posterior, because the statistical
$n_\mathrm{eff}$ is not changed by better sampling.
The width of the posterior is a property of the data, not the computer.

\begin{keybox}
\textbf{One-sentence summary:}
\emph{Statistical $n_\mathrm{eff}$ limits what you can know;
MCMC $n_\mathrm{eff}$ limits how well your computer finds it.}

\smallskip
A wide posterior for $\tau$ at $T/\tau \lesssim 2$ is not a computational failure.
It is the correct answer: the data do not contain enough information to pin $\tau$,
and the posterior is faithfully reporting that fact.
The remedy is a longer observation campaign, not a more sophisticated sampler.
\end{keybox}

% -----------------------------------------------------------------------
\subsection{Conceptual interpretation: what does a dependent observation contribute?}
\label{subsec:neff_conceptual}

The mathematics of $n_\mathrm{eff}$ and Fisher information establish that the precision
of any estimator of $\tau$ is bounded by $1/\sqrt{n_\mathrm{eff}}$, regardless of
whether the method is Bayesian posterior inference, maximum likelihood estimation,
or least squares.
This universality deserves a direct conceptual explanation.

\subsubsection*{The principle is method-independent}

The Cram\'er--Rao bound is a property of the \emph{likelihood function}, not of the
estimation method.
It says: no matter what you do with the data --- Bayesian posterior, MLE, least
squares, method of moments --- you cannot extract more precision about a parameter
than the Fisher information allows.
The method determines how close you get to the bound
(MLE achieves it asymptotically; Bayesian posteriors achieve it via
Bernstein--von Mises; ordinary least squares achieves it only for Gaussian linear
models), but the bound itself is set by the data alone.
\textbf{Every method faces the same information floor, set by $n_\mathrm{eff}$.}

\subsubsection*{Decomposing a dependent observation}

An observation $X_t$ in an AR(1) process can be split into two parts:
\begin{equation}
    X_t
    = \underbrace{\mu + \phi\,(X_{t-1} - \mu)}_{\text{predictable from past}}
    + \underbrace{V_t}_{\text{genuinely new information}}.
\end{equation}
The first part is perfectly determined by $X_{t-1}$ and the model parameters.
If you already know $X_{t-1}$ and $\phi$, this part adds \emph{zero} new information.
The second part --- the innovation $V_t \sim \mathcal{N}(0, \sigma^2)$ --- is
genuinely independent of everything that came before.
This is the new information in $X_t$.

Each time step therefore contributes exactly \textbf{one independent innovation}
to the dataset.
Over $T$ steps there are $T-1$ independent innovations, which is why
$\mathcal{I}_T(\phi) \propto T$: each innovation is one truly independent piece of
information about the conditional structure of the process.

\subsubsection*{Why $\tau$ requires more than innovations}

Knowing $\phi$ precisely is not the same as knowing $\tau$ precisely.
$\tau$ is not a property of individual innovations.
It is a property of the \textbf{rate at which the correlation between $X_t$ and
$X_{t-k}$ decays as a function of lag $k$}.
To measure a decay rate, you need to observe the process evolve through multiple
full correlation lengths: you need to see the process start somewhere, return to near
its mean, start again, and so on.

Each time the process completes one full correlation length, you obtain one
approximately independent \emph{realisation of the autocorrelation decay pattern} ---
one data point about $\tau$.
With $T = 100$ and $\tau = 38$, the process completes approximately $T/(2\tau)
\approx 1.3$ such cycles.
You have 1.3 independent observations of the decay pattern.
You cannot reliably fit an exponential decay curve to a single observation.

\subsubsection*{Which parameters benefit from which level of information}

The distinction between innovation-level and cycle-level information separates
parameters that are easy to constrain from those that are hard:

\begin{table}[h!]
\centering
\caption{Information source and scaling for each AR(1) parameter.
         Innovation-level parameters are constrained by every time step;
         timescale parameters require completed correlation cycles.}
\label{tab:info_levels}
\begin{tabular}{llll}
\toprule
Parameter & What provides information & Source & Scales as \\
\midrule
$\phi$ (AR coefficient) & Each innovation $V_t$ & 1 per step, i.i.d. & $T$ \\
$\sigma$ (innovation std) & Variance of innovations & 1 per step & $T$ \\
$\tau$ (timescale) & Each completed correlation cycle & 1 per $2\tau$ steps & $T/(2\tau) = n_\mathrm{eff}$ \\
$\mu$ (long-run mean) & Level of the process & $n_\mathrm{eff}$ per series & $n_\mathrm{eff}$ \\
\bottomrule
\end{tabular}
\end{table}

This explains the asymmetry in your simulation results: $\sigma_\tau$ is recovered
well across nearly all combinations (it uses innovation-level information, scales
with $T$), while $\tau$ itself is poorly constrained at large values (it requires
cycle-level information, scales with $n_\mathrm{eff}$).
Both are estimated from the \emph{same data}, but from \emph{different levels of
that data}.

\subsubsection*{Dependent observations are not uniformly useless}

The dependent part of $X_t$ --- the part predictable from $X_{t-1}$ --- is not
useless for \emph{all} parameters.
Its usefulness depends on what the parameter controls:

\begin{itemize}[noitemsep]
    \item For $\phi$ and $\sigma$: the dependent (predictable) part adds nothing,
          because all information about the conditional structure is already in the
          innovations.
    \item For $\mu$: the predictable part contributes, because $\mu$ determines
          where the process tends to return to.
          However, this is still limited by $n_\mathrm{eff}$, because correlated
          observations of the long-run level are redundant with each other.
    \item For $\tau$: the predictable part contributes only via the aggregate
          pattern over many cycles.
          The \emph{sequence} of predictable parts encodes the temporal correlation
          structure, but extracting the decay rate requires $n_\mathrm{eff}$ independent
          cycles.
\end{itemize}

The general principle is: \textbf{a dependent observation contributes new information
on parameter $\theta$ only to the extent that the predictable part of the observation
tells you something about $\theta$ that you did not already know.}
For parameters controlling innovation-level structure, the predictable part is
redundant and the dependent observation contributes nothing beyond the innovation.
For parameters controlling temporal scale, the dependent part contributes but only
weakly, in proportion to $n_\mathrm{eff}$.

\begin{keybox}
\textbf{Connecting $n_\mathrm{eff}$ to independence, conceptually.}

\smallskip
Each observation in a correlated time series can be split into a predictable part
(carrying no new information) and an innovation (genuinely independent).
$n_\mathrm{eff}$ counts how many of those innovations are independent \emph{for the
purpose of estimating the parameter of interest}.
For timescale parameters like $\tau$, the relevant unit of independence is not an
individual innovation but a full correlation cycle --- of which there are only
$T/(2\tau)$ in a series of length $T$.
This is why the variance of $\hat\tau$ scales as $\tau^2/n_\mathrm{eff}$ and not as
$\tau^2/T$: the parameter $\tau$ is informed by cycle-level, not
innovation-level, information.
\end{keybox}

% ──────────────────────────────────────────────────────────────────────────────
\section{General Multivariate Analysis}

The same principles extend beyond time series to any regression, hierarchical model,
or multi-parameter inference problem.

\subsection{The Fisher information matrix: the universal pre-analysis tool}

For any parametric model with parameter vector $\boldsymbol{\theta}$, the
\textbf{Cram\'er--Rao lower bound} sets the minimum variance achievable by any
unbiased estimator:

\begin{equation}
    \mathrm{Var}(\hat{\theta}_i) \geq \left[\mathcal{I}(\boldsymbol{\theta})^{-1}\right]_{ii}
\end{equation}

where $\mathcal{I}(\boldsymbol{\theta})$ is the \textbf{Fisher information matrix}
(FIM), with entries:

\begin{equation}
    \mathcal{I}_{ij}(\boldsymbol{\theta}) = -\,\mathbb{E}\!\left[
        \frac{\partial^2 \log p(\mathbf{x} \mid \boldsymbol{\theta})}{\partial \theta_i \,\partial \theta_j}
    \right]
\end{equation}

The FIM can be computed \textbf{analytically or numerically before collecting any data},
given only:
\begin{itemize}[noitemsep]
    \item the assumed generative model,
    \item a prior belief about the true parameter values,
    \item the planned data size and quality (e.g., expected photon counts, sample size).
\end{itemize}

The diagonal of $\mathcal{I}^{-1}$ gives the \emph{best-case} uncertainty on each
parameter. If this best-case exceeds your scientific precision requirement, no
inference method will achieve it --- you need more or better data.

\subsection{Parameter identifiability and degeneracies}

Before running inference, verify that each parameter is \textbf{identifiable}: that
different values of $\theta_i$ produce genuinely different observable predictions.
Non-identifiability manifests as:

\begin{itemize}[noitemsep]
    \item A near-singular FIM (very large condition number $\kappa(\mathcal{I})$).
    \item Posterior distributions that are nearly identical to the prior, regardless
          of data size.
    \item Strong correlations between parameters in the posterior.
\end{itemize}

\textbf{Practical example from X-ray spectroscopy:}
Column density $N_H$ and photon index normalization $K$ both affect spectral shape.
At low count rates, their effects are degenerate: the FIM is nearly singular in the
$N_H$--$K$ direction, meaning no finite dataset can disentangle them. This degeneracy
is detectable in advance by inspecting off-diagonal FIM elements or by computing
predicted posterior contours analytically.

\subsection{The prior-dominated versus likelihood-dominated diagnostic}

A simple pre-analysis diagnostic: compute the ratio of the \emph{expected} posterior
width to the prior width for each parameter:

\begin{equation}
    r_i = \frac{\sigma_\mathrm{posterior, \, expected}^{(i)}}{\sigma_\mathrm{prior}^{(i)}}
\end{equation}

\begin{itemize}[noitemsep]
    \item $r_i \ll 1$: the data are informative; the posterior is substantially
          narrower than the prior. \textbf{Good.}
    \item $r_i \approx 1$: the data are essentially uninformative; the posterior
          echoes the prior. The parameter is unresolvable with this dataset.
\end{itemize}

If $r_i \approx 1$ for a parameter you care about, you must either increase data
quality/quantity, impose an external constraint (tighter prior), or drop that
parameter as a free variable.


% ──────────────────────────────────────────────────────────────────────────────
\section{Pre-Analysis Feasibility Checklist}

Work through the following questions \textit{before} designing the observation campaign,
collecting data, or running any MCMC sampler.

\begin{enumerate}[label=\textbf{\arabic*.}, leftmargin=*, itemsep=6pt]

    \item \textbf{Identifiability}
    \begin{itemize}[noitemsep]
        \item Can each parameter be uniquely determined in principle, even with infinite data?
        \item Are there degeneracies that require auxiliary data or stronger priors to break?
        \item Compute or inspect off-diagonal elements of the expected FIM for large
              cross-correlations.
    \end{itemize}

    \item \textbf{Effective data volume}
    \begin{itemize}[noitemsep]
        \item For time series: compute $n_\mathrm{eff}$ via equation~\eqref{eq:neff}
              for the worst-case $\tau$ you wish to constrain. Require $n_\mathrm{eff} \gtrsim 10$.
        \item For independent observations: require $N \gg p$ (number of data points
              much larger than the number of free parameters; a factor of 10--20$\times$
              is a practical minimum).
        \item If the effective data volume is insufficient, determine how much more data
              is needed \textit{before} beginning the study.
    \end{itemize}

    \item \textbf{Signal detectability}
    \begin{itemize}[noitemsep]
        \item Compute the expected signal-to-noise ratio (SNR) for each parameter of
              interest: $\mathrm{SNR}_i \sim \Delta\hat{\theta}_i / \sigma_\mathrm{noise}$.
        \item If $\mathrm{SNR} < 3$--$5$ for a key parameter, either increase the
              sensitivity, fix that parameter to its expected value, or adopt an
              informative prior.
    \end{itemize}

    \item \textbf{Cram\'{e}r--Rao / Fisher information budget}
    \begin{itemize}[noitemsep]
        \item Evaluate $\mathcal{I}(\boldsymbol{\theta})$ at expected true parameter values.
        \item Check whether $\left[\mathcal{I}^{-1}\right]_{ii}^{1/2}$ meets the precision
              target for each parameter.
        \item If not: increase $N$ (more observations), improve quality (reduce noise),
              or relax the precision requirement.
    \end{itemize}

    \item \textbf{Prior sensitivity}
    \begin{itemize}[noitemsep]
        \item Estimate the expected ratio $r_i = \sigma_\mathrm{post} / \sigma_\mathrm{prior}$
              for each parameter.
        \item Any parameter with $r_i \gtrsim 0.7$ will be largely prior-driven.
              Be explicit about this in the analysis and interpretation.
    \end{itemize}

    \item \textbf{Computational feasibility}
    \begin{itemize}[noitemsep]
        \item Given the data volume established in steps 1--4, verify that the chosen
              sampler (HMC, NUTS, etc.) is tractable.
        \item For HMC, mixing time scales with the autocorrelation length of the
              posterior. Long-$\tau$ or degenerate posteriors require longer chains
              and more warmup steps.
        \item If the dataset size needed for statistical power makes the sampler
              intractable, consider approximate methods (variational inference,
              likelihood emulation, neural posterior estimation).
    \end{itemize}

\end{enumerate}


% ──────────────────────────────────────────────────────────────────────────────
\section{Summary}

\begin{keybox}
Before running any analysis, compute how many effectively independent, informative
constraints your data contain on each parameter of interest. If that number is
substantially below $\sim$10--20 per free parameter, the posterior will be dominated
by prior assumptions rather than data, and no computational method can remedy this.

\medskip
The tools for this pre-analysis --- the Fisher information matrix, effective sample
size formulas, and signal-to-noise estimates --- are analytical and inexpensive.
They constitute a \textbf{feasibility filter}: run them first, before committing
telescope time, observing resources, or HMC compute, and they will tell you whether
the experiment is worth running and precisely what quality and quantity of data is
required.

\medskip
This principle holds universally: for time series analysis (where $n_\mathrm{eff}$
replaces $N$), for spectral fitting (where the FIM measures how distinguishable
two spectral models are), and for any regression or hierarchical model where the
number of free parameters risks exceeding the information content of the data.
\end{keybox}

\vspace{12pt}
\noindent\rule{\linewidth}{0.4pt}
{\small
\noindent\textbf{Context.} These notes were developed in the context of Hamiltonian Monte Carlo
inference of AGN X-ray variability, modelling column density $N_H$ and normalisation $\phi$
as AR(1) stochastic processes with correlation timescale $\tau$ and amplitude $\sigma$.
The effective sample size analysis in Section~2 directly explains why $\hat{\tau}$ is
well-constrained for $\tau = 2$ steps ($n_\mathrm{eff} \approx 25$) but poorly constrained
for $\tau = 38$ or $\tau = 50$ steps ($n_\mathrm{eff} \approx 1$) when the total
time series length is $T = 100$ steps.
}

% ──────────────────────────────────────────────────────────────────────────────
\appendix
\section{Why the Spectral Normalisation $K$ Is Well-Identified Even at Low Photon Counts}
\label{app:K_snr}

This appendix explains a practically important asymmetry: while parameters such as
column density $N_H$ or correlation timescale $\tau$ can be extremely hard to constrain
with modest data, the spectral normalisation $K$ (equivalently, the log-normalisation
$\phi$) is well-identified even with a few hundred photons per spectrum.
Understanding \emph{why} guards against mistakenly applying the same low-data caution
uniformly to all parameters.

\subsection*{How $K$ enters the observed counts}

The expected photon count in energy channel $i$ is:
\begin{equation}
    \lambda_i = K \cdot f_i,
\end{equation}
where $f_i$ absorbs the instrument response, effective area, absorption, and
power-law spectral shape at channel $i$. Crucially, $K$ is a \textbf{pure
multiplicative prefactor} that scales every channel by exactly the same factor
simultaneously. It does not change the spectral shape; it only changes the amplitude.

\subsection*{The total count is a sufficient statistic for $K$}

Summing over all energy channels:
\begin{equation}
    C = \sum_i \lambda_i = K \cdot \sum_i f_i = K \cdot F_\mathrm{total},
\end{equation}
where $F_\mathrm{total}$ is a constant (determined by the spectral model and instrument
response, not by $K$). Therefore $C \propto K$ exactly, and the \textbf{total photon
count $C$ alone is a sufficient statistic for $K$}. No spectral shape information is
required.

\subsection*{The signal-to-noise argument}

Since $C$ is a sum of independent Poisson random variables, it is itself Poisson with
mean $K \cdot F_\mathrm{total}$. The standard deviation of $C$ is $\sqrt{C}$, giving
a relative uncertainty:

\begin{equation}
    \boxed{
        \frac{\sigma_K}{K} = \frac{\sigma_C}{C} = \frac{\sqrt{C}}{C} = \frac{1}{\sqrt{C}}
    }
\end{equation}

\medskip
\begin{center}
\begin{tabular}{ccc}
\toprule
Total photons $C$ & Relative uncertainty on $K$ & Practical precision \\
\midrule
$100$  & $10\%$ & Well measurable \\
$400$  & $5\%$  & Well measurable \\
$1000$ & $3\%$  & Well measurable \\
\bottomrule
\end{tabular}
\end{center}
\medskip

Even $C = 100$ photons yields a 10\% relative uncertainty on $K$, which is entirely
adequate for most scientific purposes.

\subsection*{Why the signal is unambiguous}

The word ``unambiguous'' refers to the absence of degeneracy. Consider the effect of each
spectral parameter on the observed counts:

\begin{itemize}[noitemsep]
    \item \textbf{$K$:} scales every channel by the same factor.
        Effect: uniform vertical shift in the count spectrum.
    \item \textbf{$N_H$:} appears in the exponent $\exp(-N_H \cdot \sigma(E))$,
        preferentially suppressing low-energy channels.
        Effect: energy-dependent tilt of the count spectrum, strongest at soft X-rays.
    \item \textbf{$\Gamma$:} changes the power-law slope,
        differentially redistributing counts between low and high energies.
        Effect: energy-dependent tilt of a different form.
\end{itemize}

Only $K$ moves all channels simultaneously and proportionally. Its signature cannot be
reproduced by any combination of $N_H$ and $\Gamma$, because those parameters
necessarily change the \emph{shape} while $K$ only changes the \emph{amplitude}. This
geometric separability means $K$ is identifiable from the total count alone, with no
need to resolve the spectral shape at all.

\subsection*{Fisher information for $N_H$: formal derivation and self-limiting behaviour}

Unlike $K$, the column density $N_H$ enters the spectrum nonlinearly through an
exponential absorption factor, and a Poisson-variance argument on the total count
is no longer sufficient. A full Fisher information analysis is needed, and it
reveals a fundamental asymmetry between $K$ and $N_H$.

The expected count in energy channel $i$ under the absorbed power-law model is:
\begin{equation}
    \lambda_i = K \cdot g_i \cdot \exp\!\bigl(-N_H\,\sigma(E_i)\bigr),
\end{equation}
where $g_i$ encodes the instrument response and power-law spectral shape, and
$\sigma(E_i)$ is the photoelectric absorption cross-section.
In the soft X-ray band ($E \lesssim 2\,\mathrm{keV}$), $\sigma(E) \propto E^{-3}$:
only low-energy channels carry appreciable absorption signal.

Differentiating, $\partial\lambda_i/\partial N_H = -\lambda_i\,\sigma(E_i)$.
The Fisher information contributed by Poisson channel $i$ is:
\begin{equation}
    \mathcal{I}_i(N_H)
    = \frac{1}{\lambda_i}\!\left(\frac{\partial\lambda_i}{\partial N_H}\right)^{\!2}
    = \lambda_i\,\sigma(E_i)^2.
\end{equation}
Summing over all channels:
\begin{equation}
    \boxed{
        \mathcal{I}(N_H) = \sum_i \lambda_i\,\sigma(E_i)^2
    }
    \label{eq:fisher_NH}
\end{equation}
For comparison, the same derivation applied to $K$ (using
$\partial\lambda_i/\partial K = \lambda_i/K$) gives:
\begin{equation}
    \mathcal{I}(K) = \sum_i \frac{\lambda_i}{K^2} = \frac{C}{K^2},
    \label{eq:fisher_K_formal}
\end{equation}
recovering the earlier result $\sigma_K/K = 1/\sqrt{C}$ via the Cram\'er--Rao bound.
Whereas every photon contributes equally to $\mathcal{I}(K)$, each channel's
contribution to $\mathcal{I}(N_H)$ is weighted by $\sigma(E_i)^2$: only soft-band
channels carry $N_H$ information.

This leads directly to a self-limiting behaviour. The soft-band channels with the
largest $\sigma(E_i)$ --- and hence the largest potential contribution to
$\mathcal{I}(N_H)$ in equation~\eqref{eq:fisher_NH} --- are suppressed by the same
exponential factor $\exp(-N_H\,\sigma(E_i))$ appearing in $\lambda_i$ itself.
As $N_H$ increases:
\begin{equation}
    \mathcal{I}_i(N_H) = \lambda_i\,\sigma(E_i)^2
    \;\propto\; e^{-N_H\,\sigma(E_i)}\,\sigma(E_i)^2
    \;\xrightarrow{\;N_H\to\infty\;}\; 0.
\end{equation}
The photons most informative about $N_H$ are precisely those most absorbed by it.
Hard X-ray channels avoid this self-suppression but have $\sigma(E_i)^2 \approx 0$,
contributing negligibly to $\mathcal{I}(N_H)$.
The total Fisher information therefore decreases with increasing $N_H$, causing the
Cram\'er--Rao lower bound $\sigma_{N_H} \geq 1/\sqrt{\mathcal{I}(N_H)}$ to grow.
This is why $N_H$ recovery degrades at high column densities in the simulation grid,
while $K$ and $\phi$ --- governed by $\mathcal{I}(K) = C/K^2$, which improves with
every photon --- remain well-identified throughout.


% -----------------------------------------------------------------------
\subsection*{Four levels of uncertainty in the state-space model}

To understand what the per-step $10\%$--$5\%$ precision on $K_t$ implies for recovery
of the AR(1) parameters, it is essential to separate four conceptually distinct
uncertainty levels that are easily conflated.

\begin{enumerate}[leftmargin=*, itemsep=8pt]

\item \textbf{Poisson measurement noise (Level 1).}
At each time step $t$, the total photon count $C_t$ fluctuates as
$\mathrm{Poisson}(K_t F_\mathrm{total})$, with standard deviation $\sqrt{C_t}$.
This gives a per-step relative uncertainty on $K_t$:
\begin{equation}
    \frac{\sigma_{K_t}}{K_t} = \frac{1}{\sqrt{C_t}}
    \qquad\bigl(10\%\;\text{at}\;C_t = 100;\;5\%\;\text{at}\;C_t = 400\bigr).
\end{equation}
This is irreducible instrumental noise, reducible only by increasing the photon count
(longer exposures or brighter sources).

\item \textbf{Process innovation $\sigma$ (Level 2).}
The latent process evolves according to:
\begin{equation}
    K_t = \mu + \phi\,(K_{t-1} - \mu) + \varepsilon_t,
    \qquad \varepsilon_t \sim \mathcal{N}(0,\,\sigma^2).
\end{equation}
Here $\sigma$ is the standard deviation of the random innovation at each step:
genuine stochastic variability intrinsic to the source, independent of the instrument.
$\sigma$ would persist even with infinitely many photons per step, because it
represents true physical fluctuations in $K_t$.
It is entirely distinct from Level~1 noise, and is the primary physical parameter
of interest (together with $\tau$).

\item \textbf{One-step conditional variance (Level 3).}
Given $K_{t-1}$, the conditional variance of the next state is
$\mathrm{Var}(K_t \mid K_{t-1}) = \sigma^2$,
identical to the innovation variance of Level~2.
This sets the scale of step-to-step fluctuations: even with perfect knowledge of
$K_{t-1}$ and $\phi$, the next value $K_t$ is uncertain by $\sim\!\sigma$.
The condition for the AR(1) signal to be distinguishable from photon noise is
therefore $\sigma \gg K/\sqrt{C}$ (Level~2 $\gg$ Level~1).

\item \textbf{Stationary (marginal) variance (Level 4).}
After many correlated steps, the unconditional variance of $K_t$ at stationarity is:
\begin{equation}
    \sigma_\mathrm{stat}^2 = \frac{\sigma^2}{1 - \phi^2}.
    \label{eq:sigma_stat}
\end{equation}
For $\phi$ close to 1 (large $\tau$), small per-step innovations accumulate into large
correlated excursions: $\sigma_\mathrm{stat} \gg \sigma$.
This is the amplitude of variability visible in the light curve.
For the AR(1) process to be detectable above the photon noise floor:
\begin{equation}
    \sigma_\mathrm{stat} = \frac{\sigma}{\sqrt{1-\phi^2}} \gg \frac{K}{\sqrt{C}}.
    \label{eq:detectability_cond}
\end{equation}

\end{enumerate}


% -----------------------------------------------------------------------
\subsection*{The full inference chain: from photon counts to $\tau$ and $\sigma$}

The four uncertainty levels connect in a chain that carries information from raw
photon counts all the way to the posterior distributions of $\tau$ and $\sigma$.

\medskip
\noindent\textbf{Step 1: Photons $\to$ per-step posterior on $K_t$
\emph{(Level~1 noise enters here)}.}
At each step $t$, the Fisher information $\mathcal{I}(K_t) = C_t/K_t^2$
(equation~\eqref{eq:fisher_K_formal}) sets the per-step posterior width:
$\sigma_{K_t} \approx K_t/\sqrt{C_t}$.
HMC infers all $T$ latent states $\{K_t\}$ simultaneously, with the AR(1) prior
acting as a Bayesian filter that smooths adjacent estimates using the temporal
correlation structure encoded by $\phi$ and $\sigma$.

\medskip
\noindent\textbf{Step 2: Per-step posteriors $\to$ AR(1) trajectory
\emph{(Levels~1 and~4 interact)}.}
The posterior-mean sequence $\{\hat K_t\}$ reconstructs the true AR(1) trajectory.
For this reconstruction to be faithful, the AR(1) signal (Level~4) must stand above
the photon noise (Level~1).
The relevant ratio is:
\begin{equation}
    \mathrm{SNR}_\mathrm{tracking}
    = \frac{\sigma_\mathrm{stat}}{K/\sqrt{C}}
    = \frac{\sigma}{\sqrt{1-\phi^2}}\cdot\frac{\sqrt{C}}{K}.
\end{equation}
When this ratio is large, individual step estimates $\hat K_t$ lie close to the true
values and the trajectory is reliably recovered.
When it approaches unity, measurement noise dominates and the reconstructed trajectory
carries little signal.

\medskip
\noindent\textbf{Step 3: Trajectory $\to$ $\phi$ (and hence $\tau$)
\emph{($n_\mathrm{eff}$ governs, not $C$)}.}
The autocorrelation parameter $\phi = \exp(-1/\tau)$ is estimated from the temporal
covariance structure of $\{\hat K_t\}$.
Precision is governed by the effective sample size
$n_\mathrm{eff} = T(1-\phi)/(1+\phi)$, \emph{not} by the photon count $C$.
For $\tau = 2$ and $T = 100$: $n_\mathrm{eff} \approx 25$, and $\phi$ is well-constrained.
For $\tau = 50$ and $T = 100$: $n_\mathrm{eff} \approx 1$, and $\tau$ is essentially
unconstrained regardless of how many photons are collected per step
(see Table~\ref{tab:neff_examples}).
Increasing the photon count improves the per-step trajectory reconstruction (Step~2)
but cannot substitute for insufficient temporal coverage.

\medskip
\noindent\textbf{Step 4: Trajectory and $\phi$ $\to$ $\sigma$
\emph{(Levels~1 and~2 must be separated)}.}
With $\phi$ estimated, $\sigma$ is inferred from the variance of the AR(1) residuals:
\begin{equation}
    \hat\sigma^2 \approx \frac{1}{T}\sum_{t=2}^{T}
    \bigl(K_t - \hat\mu - \hat\phi\,(K_{t-1} - \hat\mu)\bigr)^2.
\end{equation}
The precision of $\hat\sigma$ scales as $\sigma/\sqrt{n_\mathrm{eff}}$.
Beyond precision, a lower bound on detectable $\sigma$ is also set by Level~1:
if $\sigma \lesssim K/\sqrt{C}$, photon fluctuations are indistinguishable from
process variability, inflating $\hat\sigma$ above its true value.
The requirement $\sigma \gg K/\sqrt{C}$ (equation~\eqref{eq:detectability_cond}) is
therefore a necessary precondition for reliable $\sigma$ recovery.

\medskip
\noindent\textbf{Closing the chain.}
The $10\%$ per-step photon noise (100 photons) defines the noise floor that the
physical variability amplitude must clear.
Whether it does so depends on the stationary amplitude $\sigma/\sqrt{1-\phi^2}$
relative to $K/\sqrt{C}$.
Once cleared, HMC jointly infers all $T$ latent states and both AR(1) hyperparameters
in a single probabilistic model, correctly propagating Level~1 measurement noise
through to the Level~2 and Level~4 uncertainties in $\sigma$ and $\tau$.
The $10\%$--$5\%$ SNR result thus defines the \emph{minimum required physical
variability amplitude}, not merely a precision target for individual spectral fits.

\begin{keybox}
\textbf{Pre-analysis implications.}

\smallskip
\begin{itemize}[noitemsep]
    \item \textbf{$K$ vs.\ $N_H$ asymmetry:}
    All photons contribute equally to $\mathcal{I}(K) = C/K^2$, giving
    $\sigma_K/K = 1/\sqrt{C}$ at any energy.
    For $N_H$, only soft-band photons contribute to $\mathcal{I}(N_H)$, and those
    photons are self-suppressed by absorption at high column densities.
    Assess each parameter separately via its own Fisher information channel.

    \item \textbf{Four uncertainty levels:}
    Always distinguish (1)~Poisson measurement noise $\sim 1/\sqrt{C}$ per step,
    (2)~AR(1) innovation $\sigma$, (3)~one-step conditional variance $= \sigma^2$,
    and (4)~stationary variance $\sigma^2/(1-\phi^2)$.
    Reliable AR(1) recovery requires Level~1 $\ll$ Level~2; detectable variability
    requires Level~4 $\gg$ Level~1.

    \item \textbf{Noise floor vs.\ variability amplitude:}
    The $10\%$--$5\%$ per-step SNR sets the noise floor for tracking the latent
    time series.
    For AR(1) parameters to be recoverable, the stationary amplitude
    $\sigma/\sqrt{1-\phi^2}$ must substantially exceed $K/\sqrt{C}$.
    $\tau$ recovery is then further limited by effective sample size $n_\mathrm{eff}$,
    independently of photon count.
\end{itemize}
\end{keybox}

% ──────────────────────────────────────────────────────────────────────────────
\subsection*{The Cram\'er--Rao bound: what it is and what it requires}
\label{sec:cramer_rao}

\paragraph{What it is.}
The \textbf{Cram\'er--Rao Lower Bound} (CRLB) is a \emph{theorem} in mathematical
statistics, proved independently by C.R.~Rao (1945) and Harald~Cram\'er (1946).
It states that for any unbiased estimator $\hat\theta$ of a parameter $\theta$,

\begin{equation}
    \mathrm{Var}(\hat\theta) \;\geq\; \frac{1}{\mathcal{I}_T(\theta)},
    \label{eq:crlb}
\end{equation}

where $\mathcal{I}_T(\theta)$ is the total Fisher information in the sample.
The bound is an \emph{inequality}: it tells you the lowest variance that any
estimator in the admissible class can achieve, not the variance of a particular
estimator you have in hand.
An estimator that achieves equality in~\eqref{eq:crlb} is called \textbf{efficient}.
The maximum-likelihood estimator (MLE) is asymptotically efficient: it reaches
the CRLB as $T\to\infty$.

For a \emph{biased} estimator with bias $b(\theta) = E[\hat\theta] - \theta$,
the generalised form is:
\begin{equation}
    \mathrm{Var}(\hat\theta) \;\geq\;
    \frac{\bigl[1 + db(\theta)/d\theta\bigr]^2}{\mathcal{I}_T(\theta)}.
\end{equation}

\paragraph{Regularity conditions.}
The CRLB does \emph{not} apply to every estimator or every model automatically.
Two separate sets of conditions must hold.

\medskip
\noindent\textit{(A) Conditions on the estimator itself:}

\begin{enumerate}[noitemsep]
    \item \textbf{Unbiasedness.} $E_\theta[\hat\theta] = \theta$ for all $\theta$
    in the parameter space.
    This is the primary condition; it defines the class of estimators to which the
    bound applies.
    \item \textbf{Finite variance.} $\mathrm{Var}(\hat\theta) < \infty$.
\end{enumerate}

\medskip
\noindent\textit{(B) Regularity conditions on the statistical model $f(x;\theta)$:}

\begin{table}[h]
\centering
\caption{Regularity conditions required for the Cram\'er--Rao bound to hold.
If any of these fails, the bound may not exist or may be violated by some
estimator, and an alternative bound (Table~\ref{tab:bounds_family}) should
be used instead.}
\label{tab:regularity}
\renewcommand{\arraystretch}{1.4}
\begin{tabular}{p{4.2cm} p{9.0cm}}
\hline
\textbf{Condition} & \textbf{Meaning} \\
\hline
Open parameter space &
$\theta \in \Theta$ where $\Theta \subseteq \mathbb{R}$ is an open set;
the bound is a local (differential) statement. \\
Support independence &
The set $\{x : f(x;\theta)>0\}$ does not depend on $\theta$.
If the support moves with $\theta$ (e.g.\ Uniform$(0,\theta)$),
the CRLB does not hold. \\
Differentiability &
$\ln f(x;\theta)$ is differentiable with respect to $\theta$ for
almost all $x$; the score function
$s(x;\theta) = \partial\ln f/\partial\theta$ exists. \\
Interchange of $\int$ and $\partial/\partial\theta$ &
Differentiation through the integral
$\partial/\partial\theta\!\int\! f(x;\theta)\,dx = 0$
is valid (dominated convergence); this ensures
$E_\theta[s(x;\theta)] = 0$. \\
Finite, positive Fisher information &
$0 < \mathcal{I}(\theta) = E_\theta[s(x;\theta)^2] < \infty$. \\
\hline
\end{tabular}
\end{table}

\medskip
\noindent\textit{What regularity conditions mean conceptually.}
The CRLB is derived by differentiating an integral identity with respect to $\theta$.
For this operation to be valid, the model must be \emph{smooth} in $\theta$: the
distribution must change in a well-defined, differentiable way as the parameter
changes, and that change must be measurable through the score function.
When the support of $f$ depends on $\theta$ --- for example, a uniform distribution
on $[0,\theta]$ --- moving $\theta$ \emph{changes which data values are possible},
not just how probable they are.
In this case, the data carry information about $\theta$ through a completely
different mechanism (extreme-order statistics), and the Fisher-information route
breaks down.
Similarly, if $f$ is not differentiable in $\theta$ (e.g.\ a double-exponential
location model at the point of non-differentiability), the score function does not
exist and the CRLB cannot be formed.

\paragraph{How much memory survives the observation window?}
A concrete consequence of autocorrelation is that the process retains memory
across the entire observation window $T$.
The autocorrelation at lag $k$ for an AR(1) process is
$\rho(k) = \exp(-k/\tau)$.
Evaluating this at $k = T = 100$ for each value of $\tau$ in the prior support
shows how much of the original correlation survives after all $T$ steps:

\begin{table}[h]
\centering
\caption{Fraction of the original autocorrelation remaining after the full
$T = 100$ step observation window, $\rho(T) = \exp(-T/\tau)$, for selected
values of $\tau$.
For large $\tau$ the memory is far from exhausted at the end of the window,
which is the direct cause of the limited effective sample size.
Note that $\tau = 80$ gives $\rho(100) \approx 0.29$: nearly 30\% of the
original correlation survives after 100 steps.}
\label{tab:memory_survival}
\renewcommand{\arraystretch}{1.4}
\begin{tabular}{cccc}
\hline
$\tau$ & $\varphi = e^{-1/\tau}$ &
$\rho(100) = e^{-100/\tau}$ &
Memory remaining \\
\hline
2  & 0.607 & $e^{-50} \approx 0$        & Effectively zero after $\sim$6 steps \\
10 & 0.905 & $e^{-10} \approx 0.000045$ & Gone \\
40 & 0.975 & $e^{-2.5} \approx 0.082$   & 8\% remains \\
80 & 0.988 & $e^{-1.25} \approx 0.287$  & \textbf{29\% remains} \\
\hline
\end{tabular}
\end{table}

\noindent For $\tau = 80$ (the upper limit of the prior), 29\% of the original
correlation is still present at step 100.
This is why the process at large $\tau$ is described as ``near a random walk'':
$\varphi = e^{-1/80} \approx 0.988$ is only 1.2\% below unity, and the
correlation decays by less than $e^{-1/80} \approx 1.2\%$ \emph{multiplicatively}
(not additively) per step.
After 100 steps the accumulated decay is $0.988^{100} = e^{-100/80} \approx 0.29$,
leaving substantial memory intact.

Note that \textbf{consistency} is \emph{not} a requirement for the CRLB.
An estimator can be consistent or non-consistent; the bound applies equally to both,
as long as unbiasedness and the regularity conditions hold.
Consistency and efficiency are independent properties: the MLE is asymptotically
consistent \emph{and} asymptotically efficient, but these follow from separate
arguments.

% ──────────────────────────────────────────────────────────────────────────────
\subsection*{A family of lower bounds: pre-analysis tools}
\label{sec:bounds_family}

The CRLB is the most widely used lower bound but not the only one.
Several related bounds address its limitations --- most importantly, the fact
that the CRLB is often not tight (an efficient estimator may not exist) and
that it fails when regularity conditions are violated.
All bounds in this family answer the same \emph{pre-analysis question}:

\begin{center}
\emph{Before running any inference, what is the minimum variance that any
estimator can achieve, given this data and this model?}
\end{center}

\noindent If the bound itself is large, no inference method --- frequentist or
Bayesian, MLE or MCMC or SBI --- can do better.
This separates \emph{fundamental statistical limits} from
\emph{suboptimality of the chosen estimator}.

\begin{table}[h]
\centering
\caption{Family of lower bounds on estimator variance, ordered from least to
most tight (within the frequentist class) and including the Bayesian extension.
All bounds require finite Fisher information to exist; the Chapman--Robbins bound
is the exception and applies even when regularity conditions fail.}
\label{tab:bounds_family}
\renewcommand{\arraystretch}{1.5}
\begin{tabular}{p{3.8cm} p{1.6cm} p{4.8cm} p{3.2cm}}
\hline
\textbf{Bound} & \textbf{Type} & \textbf{Key idea} & \textbf{When to use} \\
\hline
Cram\'er--Rao (1945/46) &
Freq. &
$1/\mathcal{I}_T(\theta)$; first-order score function. &
Standard starting point; requires regularity. \\
Bhattacharyya (1946) &
Freq. &
Uses higher-order derivatives of log-likelihood; strictly tighter than CRLB. &
When CRLB is not achieved by any estimator. \\
Chapman--Robbins (1951) &
Freq. &
Finite-difference version; no differentiability needed. &
When regularity conditions fail (e.g.\ uniform support). \\
Barankin (1949) &
Freq. &
Tightest classical bound; least-favourable parametric submodel. &
When sharpest possible floor is required. \\
Van~Trees inequality (1968) &
Bayesian &
Adds prior: $1/(\mathcal{I}_T(\theta)+\mathcal{I}_\mathrm{prior}(\theta))$. &
Bayesian setting; directly bounds posterior variance. \\
Ziv--Zakai (1969) &
Bayesian &
Links estimation error to hypothesis-testing error probability. &
Signal detection and communications problems. \\
Weiss--Weinstein (1988) &
Bayesian &
Tighter than Van~Trees in most practical cases. &
When Van~Trees bound is too loose. \\
\hline
\end{tabular}
\end{table}

\paragraph{Which bound is most relevant here.}
The \textbf{Van~Trees inequality} is the most directly applicable bound for the
inference in this project.
Because the prior on $\tau$ is log-uniform on $[2, 80]$, it contributes prior
information $\mathcal{I}_\mathrm{prior}(\tau)$ that slightly tightens the bound
beyond the data alone.
The Bayesian bound reads:
\begin{equation}
    E_\theta\bigl[\mathrm{Var}(\hat\tau \mid \mathbf{X})\bigr]
    \;\geq\;
    \frac{1}{\mathcal{I}_T(\tau) + \mathcal{I}_\mathrm{prior}(\tau)}.
\end{equation}
However, when $n_\mathrm{eff} \approx 1.3$ (as for $\tau=38$, $T=100$),
$\mathcal{I}_T(\tau) = T/(2\tau^3) \approx 0.023$ is so small that the prior
term dominates and the posterior width is controlled almost entirely by the
prior support, not by the data.
This is the formal reason why the posterior for $\tau_{NH}$ at $\tau=38$ is wide
regardless of photon count or inference algorithm.

\begin{keybox}
\textbf{Pre-analysis interpretation of the bound family.}

\smallskip
\begin{itemize}[noitemsep]
    \item \textbf{The bounds are not about a specific method.}
    They apply to every possible unbiased (or, in the Bayesian case, any) estimator.
    A wide bound means the problem is hard; no method can escape it.
    \item \textbf{Check regularity first.}
    Before applying the CRLB, verify that the support of the model does not depend
    on the parameter, that the log-likelihood is differentiable, and that Fisher
    information is finite.
    For the AR(1) model used here, all regularity conditions hold for
    $\phi \in (-1,1)$ and $\sigma > 0$.
    \item \textbf{Use the Van~Trees bound for Bayesian inference.}
    When a prior is available, the Van~Trees inequality is the correct pre-analysis
    tool; the frequentist CRLB ignores prior information and can be overly
    pessimistic.
    \item \textbf{Consistency $\neq$ efficiency.}
    An estimator can converge to the truth (consistent) without achieving the CRLB
    (efficient).
    The MLE is both, but only asymptotically as $T\to\infty$.
    For finite $T$ --- as in all practical time series --- a gap between the MLE
    variance and the CRLB typically remains.
\end{itemize}
\end{keybox}

% ──────────────────────────────────────────────────────────────────────────────
\subsection*{Pre-analysis diagnostic~1: identifiability analysis}
\label{sec:identifiability}

Identifiability is logically prior to everything else: precision is irrelevant if
parameters cannot be uniquely recovered even with infinite data.
There are two levels.

\medskip
\noindent\textbf{Structural identifiability.}
The model $f(x;\theta)$ is \emph{structurally identifiable} if the map
$\theta \mapsto f(\cdot;\theta)$ is injective: distinct parameter values produce
distinct data distributions.
If two parameter vectors $\theta \neq \theta'$ satisfy
$f(x;\theta) = f(x;\theta')$ for all $x$, the parameters are fundamentally
confounded and no amount of data or algorithmic sophistication can separate them.
This is a property of the model, not of the data.

\medskip
\noindent\textbf{Practical (local) identifiability.}
Even when structural identifiability holds, a parameter may be practically
unidentifiable for a given dataset: the likelihood surface is so flat in some
direction that no estimator can pin it down.
The Fisher information matrix provides the operational test:

\begin{itemize}[noitemsep]
    \item \textbf{FIM rank test.}
    If $\det\mathcal{I}(\theta) = 0$, at least one parameter direction carries
    exactly zero Fisher information --- exact non-identifiability.
    \item \textbf{Condition number.}
    $\kappa = \lambda_{\max}/\lambda_{\min}$ of $\mathcal{I}(\theta)$.
    Large $\kappa$ signals near-non-identifiability: the bound exists but is
    enormous along the ill-conditioned eigenvector direction.
    A rule of thumb from numerical linear algebra: $\kappa > 10^3$ indicates
    severe ill-conditioning.
    \item \textbf{Profile likelihood.}
    Fix one parameter and maximise the likelihood over the rest.
    A flat profile indicates a non-identifiable direction that the FIM rank
    test may miss in practice due to finite-sample effects.
    \item \textbf{Correlation matrix of estimators.}
    $\mathrm{Corr}(\hat\theta_i,\hat\theta_j) =
    -[\mathcal{I}^{-1}]_{ij}/\sqrt{[\mathcal{I}^{-1}]_{ii}[\mathcal{I}^{-1}]_{jj}}$.
    Values close to $\pm 1$ reveal parameter trade-offs: the data constrain a
    \emph{combination} of the two parameters well, but not each individually.
\end{itemize}

\noindent\textit{Relevance here.}
For the AGN model, structural identifiability holds: $N_H$ and $K$ enter the
spectrum through physically distinct mechanisms (exponential absorption vs.\ power-law
normalisation) and at different energies.
However the FIM is nearly block-diagonal in energy bands, so the off-diagonal
$\mathcal{I}(N_H, K)$ element is small but non-zero, meaning a mild correlation
between $\hat N_H$ and $\hat K$ is expected.

% ──────────────────────────────────────────────────────────────────────────────
\subsection*{Pre-analysis diagnostic~2: information-theoretic minimax bounds}
\label{sec:minimax}

The CRLB family bounds the variance of a specific class of estimators (unbiased,
or Bayesian).
The \emph{minimax} perspective asks a harder question: what is the best
\emph{worst-case} error any estimator can achieve, over all possible parameter values?
These bounds are completely method-agnostic and require no regularity conditions.

\paragraph{Mutual information as an absolute ceiling.}
The mutual information $I(\theta; \mathbf{X})$ between the parameter and the data
is the maximum amount of information about $\theta$ that the data can ever provide,
regardless of the estimator or the inference method.
By the data processing inequality:
\begin{equation}
    I(\theta;\, g(\mathbf{X})) \leq I(\theta;\, \mathbf{X})
    \quad \text{for any function } g.
\end{equation}
Any summary statistic, dimensionality reduction, or compression of the data can
only \emph{lose} information, never gain it.
Computing $I(\theta;\mathbf{X})$ before inference therefore sets an absolute
ceiling on what can be learned.

\paragraph{Fano's inequality.}
For a discretised parameter taking $M$ values, Fano's inequality gives a lower
bound on the probability of error $P_e$ for any estimator:
\begin{equation}
    P_e \;\geq\; 1 - \frac{I(\theta;\mathbf{X}) + \log 2}{\log M}.
\end{equation}
This is particularly useful for detection problems: if the mutual information is
small relative to $\log M$, the error rate is bounded away from zero regardless
of the inference algorithm used.

\paragraph{Le~Cam's method and Assouad's lemma.}
These provide minimax lower bounds by reducing estimation to hypothesis testing:

\begin{itemize}[noitemsep]
    \item \textbf{Le~Cam's method.}
    Find the two parameter values $\theta_0$, $\theta_1$ that are hardest to
    distinguish (smallest total variation distance between $f(\cdot;\theta_0)$
    and $f(\cdot;\theta_1)$).
    The minimax risk is bounded below by the error rate of the best two-point test.
    \item \textbf{Assouad's lemma.}
    Multivariate generalisation: constructs a $2^d$ hypercube of parameter values
    and bounds the minimax risk in terms of the pairwise distinguishability.
    Gives the sharpest known minimax lower bounds for many problems.
\end{itemize}

\noindent These bounds are strictly stronger than the CRLB in irregular models
(where the CRLB does not apply) and provide non-asymptotic guarantees.

% ──────────────────────────────────────────────────────────────────────────────
\subsection*{Pre-analysis diagnostic~3: sensitivity and conditioning}
\label{sec:sensitivity}

Even when parameters are identifiable and the CRLB is finite, the inference
problem can be numerically and statistically ill-conditioned: small perturbations
in the data produce large swings in the estimates.
These diagnostics quantify that amplification.

\begin{table}[h]
\centering
\caption{Sensitivity and conditioning diagnostics. All are computable from the
Fisher information matrix before any data is collected.}
\label{tab:sensitivity}
\renewcommand{\arraystretch}{1.4}
\begin{tabular}{p{4.0cm} p{9.2cm}}
\hline
\textbf{Diagnostic} & \textbf{What it reveals} \\
\hline
FIM eigendecomposition &
Eigenvectors identify which linear combinations of parameters are
well-constrained (large eigenvalue) vs.\ poorly constrained (small eigenvalue).
Eigenvalues give the variance in each principal direction. \\
Condition number $\kappa = \lambda_{\max}/\lambda_{\min}$ &
Ratio of best to worst constrained direction; $\kappa \gg 1$ means
the problem is ill-conditioned and small noise in data is amplified into
large uncertainty in the hard direction. \\
Estimator correlation matrix &
Off-diagonal entries $\mathrm{Corr}(\hat\theta_i,\hat\theta_j)$ near
$\pm 1$ reveal parameter degeneracies: the data constrain a combination
but not each individually. \\
Variance inflation factor (VIF) &
In regression: $\mathrm{VIF}_j = 1/(1-R_j^2)$ where $R_j^2$ is the
$R^2$ from regressing predictor $j$ on the others.
$\mathrm{VIF} > 10$ indicates harmful collinearity. \\
Sensitivity matrix $\partial E[\mathbf{X}]/\partial\theta$ &
Directly measures how much the expected data changes per unit change in
the parameter; flat sensitivity = hard to identify. \\
\hline
\end{tabular}
\end{table}

\noindent\textit{Relevance here.}
The FIM eigendecomposition for the joint $(N_H, K)$ model reveals whether there
is a ``soft direction'' --- a combination of $N_H$ and $K$ that is poorly
constrained at a given count level.
At low counts, the two parameters become more correlated because both $N_H$
(through absorption) and $K$ (through normalisation) can explain a reduction
in observed counts.

% ──────────────────────────────────────────────────────────────────────────────
\subsection*{Pre-analysis diagnostic~4: sample size and power analysis}
\label{sec:power}

Once the bound on variance is known, one can ask: how much data is needed to
achieve a desired precision or detection probability?

\paragraph{Classical power analysis.}
Given an effect size $\delta$, significance level $\alpha$, and desired power
$1-\beta$, compute the minimum sample size $T$ such that a test of size $\alpha$
achieves power $1-\beta$.
This is standard in hypothesis-testing contexts (e.g.\ detecting variability
vs.\ no variability).

\paragraph{Bayesian sample size determination.}
The Bayesian analogue targets the \emph{expected posterior variance}:
find the minimum $T$ such that
\begin{equation}
    E_{\mathbf{X}}\bigl[\mathrm{Var}(\theta \mid \mathbf{X})\bigr]
    \;\leq\; \epsilon^2,
\end{equation}
where $\epsilon$ is the desired posterior standard deviation.
Using the Van~Trees bound as an approximation, this gives an analytic target
before any data are collected.

\paragraph{$n_\mathrm{eff}$ requirement for autocorrelated processes.}
This is the most directly actionable tool for the AR(1) setting.
The Cram\'er--Rao bound gives $\mathrm{std}(\hat\tau)/\tau \geq 1/\sqrt{n_\mathrm{eff}}$,
so achieving relative precision $\delta$ requires:
\begin{equation}
    n_\mathrm{eff} = \frac{T}{2\tau} \;\geq\; \frac{1}{\delta^2}
    \quad\Longrightarrow\quad
    T \;\geq\; \frac{2\tau}{\delta^2}.
    \label{eq:T_requirement}
\end{equation}
Table~\ref{tab:T_requirement} evaluates this for the parameter combinations
studied in this project.

\begin{table}[h]
\centering
\caption{Minimum time series length $T_{\min} = 2\tau/\delta^2$ required to
achieve relative precision $\delta = \mathrm{std}(\hat\tau)/\tau$ at the
Cram\'er--Rao limit, for selected values of $\tau$ and $\delta$.
The current dataset has $T = 100$.}
\label{tab:T_requirement}
\renewcommand{\arraystretch}{1.3}
\begin{tabular}{cccccc}
\hline
 & \multicolumn{5}{c}{Required $T_{\min}$ for target relative precision $\delta$} \\
\cmidrule(lr){2-6}
$\tau$ & $\delta=0.5$ & $\delta=0.3$ & $\delta=0.2$ & $\delta=0.1$ & $\delta=0.05$ \\
\hline
2   &  16  &  44  &  100 &  400  &  1\,600 \\
10  &  80  &  222 &  500 &  2\,000 & 8\,000 \\
38  &  304 &  844 & 1\,900 & 7\,600 & 30\,400 \\
50  &  400 & 1\,111 & 2\,500 & 10\,000 & 40\,000 \\
80  &  640 & 1\,778 & 4\,000 & 16\,000 & 64\,000 \\
\hline
\end{tabular}
\end{table}

\noindent For $\tau=38$ and $T=100$, even the loosest target ($\delta=0.5$,
factor-of-two precision) requires $T_{\min} = 304 > 100$.
With $T=100$ the achievable precision floor is
$\delta \geq \sqrt{2\tau/T} = \sqrt{76/100} \approx 0.87$,
meaning the best any estimator can do is a relative standard deviation of $87\%$.

% ──────────────────────────────────────────────────────────────────────────────
\subsection*{Pre-analysis diagnostic~5: local asymptotic normality and model
regularity}
\label{sec:lan}

\paragraph{What is local asymptotic normality (LAN)?}
A parametric model satisfies \textbf{LAN} (Le~Cam 1960) if, locally around any
$\theta_0$, the log-likelihood ratio for a $1/\sqrt{T}$-neighbourhood of
$\theta_0$ converges in distribution to a Gaussian shift:
\begin{equation}
    \log \frac{L(\theta_0 + h/\sqrt{T})}{L(\theta_0)}
    \;\xrightarrow{d}\;
    h\,\Delta - \tfrac{1}{2}h^2 \mathcal{I}(\theta_0),
    \quad \Delta \sim \mathcal{N}(0, \mathcal{I}(\theta_0)).
\end{equation}
LAN is the condition under which the model is \emph{regular} in the asymptotic
sense.

\paragraph{Consequences when LAN holds.}
\begin{itemize}[noitemsep]
    \item All regular estimators converge at rate $1/\sqrt{T}$.
    \item The MLE is asymptotically efficient: it achieves the CRLB.
    \item The posterior is asymptotically normal (Bernstein--von Mises theorem).
    \item The CRLB is sharp: an efficient estimator exists asymptotically.
\end{itemize}

\paragraph{Consequences when LAN fails.}
In non-regular models --- change-point problems, models with non-smooth
likelihoods, or boundary parameters --- the $1/\sqrt{T}$ rate breaks down.
Convergence may be faster (e.g.\ $1/T$ for change-point location) or slower,
and the CRLB is no longer the correct floor.
Pre-analysis should verify LAN before relying on the CRLB as the fundamental
limit.

\paragraph{LAN for the AR(1) model.}
The Gaussian AR(1) model with $\phi \in (-1,1)$ and $\sigma > 0$ satisfies LAN
at every interior point of the parameter space.
The regularity conditions in Table~\ref{tab:regularity} are precisely the
sufficient conditions for LAN to hold.
At the boundary $\phi = 1$ (unit root), LAN fails: the parameter estimate of
$\phi$ converges faster than $1/\sqrt{T}$ (Dickey--Fuller theory), and the
CRLB no longer applies.
This boundary is outside our prior support ($\phi < \exp(-1/80) < 1$), so LAN
holds throughout the prior.

% ──────────────────────────────────────────────────────────────────────────────
\subsection*{Pre-analysis diagnostic~6: prior and model consistency checks}
\label{sec:prior_checks}

The diagnostics above assume the model is correctly specified.
Before inference, one should verify that the model and prior are consistent with
the data-generating process.

\paragraph{Prior predictive check.}
Simulate synthetic data from the generative model using parameters drawn from
the prior:
\begin{equation}
    \tilde\theta \sim \pi(\theta), \quad
    \tilde{\mathbf{X}} \sim f(\mathbf{x}\mid\tilde\theta).
\end{equation}
If the resulting synthetic datasets look qualitatively different from real
observations (wrong count ranges, wrong spectral shape, wrong variability
amplitude), the prior or likelihood is misspecified \emph{before} inference
starts.
This is a zero-cost diagnostic that costs only simulation time.

\paragraph{Prior--data conflict.}
After collecting data, check whether the observed $\mathbf{X}$ falls in a
low-probability region of the prior predictive distribution:
\begin{equation}
    p(\mathbf{X}) = \int f(\mathbf{X}\mid\theta)\,\pi(\theta)\,d\theta.
\end{equation}
A very low $p(\mathbf{X})$ signals that the prior and the data are in conflict
--- either the prior is misspecified, or the data contain unexpected features.

\paragraph{Kullback--Leibler support condition (Schwartz theorem).}
A necessary condition for Bayesian posterior consistency is that the prior
assigns positive probability to every KL-neighbourhood of the true parameter:
\begin{equation}
    \Pi\bigl(\theta : \mathrm{KL}(f(\cdot;\theta_0)\,\|\,f(\cdot;\theta)) < \varepsilon\bigr) > 0
    \quad \forall\,\varepsilon > 0.
\end{equation}
If the true $\theta_0$ lies in a region of zero prior probability (e.g.\ $\tau_0$
outside the prior support $[2,80]$), the posterior will \emph{never} concentrate
near the truth regardless of how much data are collected.
Always verify that the true parameter values used in simulation lie in the
interior of the prior support.

\paragraph{Prior sensitivity.}
Perturb the prior slightly and measure how much the posterior changes.
Large sensitivity means the data are not informative enough to overwhelm the
prior; the posterior width is prior-dominated.
This is directly related to the Van~Trees bound: when
$\mathcal{I}_T(\theta) \ll \mathcal{I}_\mathrm{prior}(\theta)$, the posterior is
prior-dominated and prior sensitivity will be high.

\begin{keybox}
\textbf{Summary: six pre-analysis diagnostics and what each rules out.}

\smallskip
\begin{tabular}{p{3.8cm} p{9.0cm}}
\textbf{Diagnostic} & \textbf{Fundamental limit revealed} \\
\hline
1.\ Identifiability (FIM rank, profile likelihood) &
Whether parameter recovery is possible at all \\
2.\ Minimax bounds (Fano, Le~Cam, mutual information) &
Best possible error for any estimator, no regularity needed \\
3.\ Sensitivity and conditioning (FIM eigenstructure, VIF) &
Which parameter directions are hard vs.\ easy \\
4.\ Sample size / $n_\mathrm{eff}$ requirement &
Minimum data needed to achieve a target precision \\
5.\ LAN / model regularity check &
Whether the CRLB is achievable asymptotically \\
6.\ Prior predictive / KL support check &
Whether the model and prior are consistent with the data \\
\end{tabular}
\end{keybox}

% ──────────────────────────────────────────────────────────────────────────────
\section{Experiment Design from First Principles}
\label{sec:experiment_design}

\textit{This section is reserved for a detailed treatment of how the pre-analysis
diagnostics above can be used to design an observational programme.
The central idea is to translate fundamental statistical limits --- the CRLB,
$n_\mathrm{eff}$, minimax bounds, and identifiability conditions --- into concrete
observational requirements: how many time steps $T$, how many photons per step $C$,
what cadence, and what source properties are needed to recover a given parameter
with a given precision.
The following topics will be developed here.}

\medskip
\begin{enumerate}[noitemsep]
    \item \textbf{From precision target to $T_{\min}$.}
    Use equation~\eqref{eq:T_requirement} to convert a desired relative precision
    $\delta$ and timescale $\tau$ into a minimum observation length.
    \item \textbf{From photon budget to per-step SNR.}
    Use the Fisher information for $N_H$ and $K$ to determine the minimum count
    rate $C$ needed so that Level-1 noise (Poisson) is below Level-2 noise
    (AR(1) innovation), enabling trajectory tracking.
    \item \textbf{Joint design: $T$ and $C$ as coupled requirements.}
    Construct iso-precision contours in the $(T, C)$ plane using the four-level
    uncertainty decomposition (Section~2).
    \item \textbf{Identifiability check at design stage.}
    Evaluate the FIM condition number as a function of $N_H$, $K$, and $C$;
    identify the regime where $N_H$ and $K$ become degenerate.
    \item \textbf{Bayesian experimental design (optimal design).}
    Maximise the expected information gain $E[D_\mathrm{KL}(\pi_\mathrm{post}\,\|\,\pi_\mathrm{prior})]$
    over the design space $(T, C, \text{cadence})$; this is the decision-theoretic
    version of the pre-analysis pipeline.
    \item \textbf{Robustness to model misspecification.}
    Assess how design choices change when the true process deviates from the
    assumed AR(1) model (e.g.\ non-stationarity, heavier tails, or a second
    timescale).
\end{enumerate}

% ──────────────────────────────────────────────────────────────────────────────
\section*{Key References}

\begin{description}[leftmargin=0pt, itemsep=4pt]

\item[Bartlett (1946)]
Bartlett, M.S.\ (1946).
On the theoretical specification of sampling properties of autocorrelated time
series.
\textit{Journal of the Royal Statistical Society, Series B}, \textbf{8}(1), 27--41.\\
\textit{Original derivation of $\mathrm{Var}(\bar{X})$ with the triangular taper
(equation~\eqref{eq:var_xbar_exact}).}

\item[Kish (1965)]
Kish, L.\ (1965).
\textit{Survey Sampling}.
Wiley, New York.\\
\textit{Original introduction of the term ``effective sample size'' $n_\mathrm{eff}$
for correlated (clustered) observations in survey statistics.}

\item[Geyer (1992)]
Geyer, C.J.\ (1992).
Practical Markov Chain Monte Carlo.
\textit{Statistical Science}, \textbf{7}(4), 473--511.\\
\textit{Defines $n_\mathrm{eff} = T/(2\tau_\mathrm{int})$ for MCMC chains and
establishes the integrated autocorrelation time as the key convergence diagnostic.
Canonical reference for both statistical and MCMC $n_\mathrm{eff}$.}

\item[Rao (1945)]
Rao, C.R.\ (1945).
Information and the accuracy attainable in the estimation of statistical parameters.
\textit{Bulletin of the Calcutta Mathematical Society}, \textbf{37}, 81--91.\\
\textit{Original proof of the lower bound on estimator variance now known as
the Cram\'er--Rao bound. Introduces Fisher information as the curvature of the
log-likelihood.}

\item[Cram\'er (1946)]
Cram\'er, H.\ (1946).
\textit{Mathematical Methods of Statistics}.
Princeton University Press.\\
\textit{Independent derivation of the same lower bound (Chapter 32). Also contains
the Bernstein--von Mises theorem and asymptotic theory of the MLE.}

\item[Bhattacharyya (1946)]
Bhattacharyya, A.\ (1946--47).
On some analogues of the amount of information and their use in statistical
estimation.
\textit{Sankhya}, \textbf{8}, 1--14 and 201--218.\\
\textit{Generalisation of the CRLB to higher-order derivatives of the
log-likelihood, giving a sequence of tighter lower bounds.}

\item[Le~Cam (1960)]
Le~Cam, L.\ (1960).
Locally asymptotically normal families of distributions.
\textit{University of California Publications in Statistics}, \textbf{3}, 37--98.\\
\textit{Introduces local asymptotic normality (LAN): the condition under which a
parametric model is locally equivalent to a Gaussian shift experiment.
Establishes that the CRLB is sharp (achievable) for all regular models.}

\item[Chapman \& Robbins (1951)]
Chapman, D.G.\ \& Robbins, H.\ (1951).
Minimum variance estimation without regularity assumptions.
\textit{Annals of Mathematical Statistics}, \textbf{22}(4), 581--586.\\
\textit{Finite-difference lower bound that does not require differentiability of
the log-likelihood; applies when the support of the distribution depends on
the parameter.}

\item[Cover \& Thomas (2006)]
Cover, T.M.\ \& Thomas, J.A.\ (2006).
\textit{Elements of Information Theory}, 2nd edition.
Wiley, New York.\\
\textit{Definitive reference for mutual information, Fano's inequality, the data
processing inequality, and channel capacity (Chapters 2, 7, and 8).
The data processing inequality $I(\theta;g(\mathbf{X})) \leq I(\theta;\mathbf{X})$
is proved in Theorem~2.8.1.}

\item[Van~Trees (1968)]
Van~Trees, H.L.\ (1968).
\textit{Detection, Estimation, and Modulation Theory, Part~I}.
Wiley, New York.\\
\textit{Introduces the Bayesian Cram\'er--Rao bound (Van~Trees inequality),
which incorporates prior information through the prior Fisher information
$\mathcal{I}_\mathrm{prior}(\theta)$. The directly applicable bound for
Bayesian posterior variance.}

\item[Casella \& Berger (2002)]
Casella, G.\ \& Berger, R.L.\ (2002).
\textit{Statistical Inference}, 2nd edition.
Duxbury Press.\\
\textit{Standard graduate text for Fisher information, the Cram\'er--Rao bound,
regularity conditions, and MLE asymptotic efficiency (Chapters 7 and 10).}

\item[Brockwell \& Davis (2002)]
Brockwell, P.J.\ \& Davis, R.A.\ (2002).
\textit{Introduction to Time Series and Forecasting}, 2nd edition.
Springer.\\
\textit{AR(1) process, autocovariance function, spectral density, and the
asymptotic distribution of the sample mean (expressed as $2\pi f(0)/T$,
equivalent to our $2\sigma^2\tau_\mathrm{int}/T$).
Does not use the term $n_\mathrm{eff}$.}

\item[Hamilton (1994)]
Hamilton, J.D.\ (1994).
\textit{Time Series Analysis}.
Princeton University Press.\\
\textit{Comprehensive treatment of AR processes, spectral theory, and estimation.
Chapters 5 and 8 cover the autocovariance-generating function and sample mean
variance in depth.}

\item[Gelman et al.\ (2013)]
Gelman, A., Carlin, J.B., Stern, H.S., Dunson, D.B., Vehtari, A.\ \& Rubin, D.B.\
(2013).
\textit{Bayesian Data Analysis}, 3rd edition.
CRC Press.\\
\textit{Covers Bayesian inference, HMC/NUTS, MCMC $n_\mathrm{eff}$, R-hat
diagnostics, and the normal approximation to the posterior
(Bernstein--von Mises, \S4.4).
The most directly useful single book for the inference methodology used here.}

\item[Brooks et al.\ (2011)]
Brooks, S., Gelman, A., Jones, G.L.\ \& Meng, X.-L.\ (eds.)\ (2011).
\textit{Handbook of Markov Chain Monte Carlo}.
CRC Press.\\
\textit{Contains Neal's definitive chapter on HMC (Chapter 5) and Geyer's
chapter on MCMC convergence theory. Reference for HMC leapfrog geometry and
MCMC $n_\mathrm{eff}$ formula.}

\item[van der Vaart (1998)]
van der Vaart, A.W.\ (1998).
\textit{Asymptotic Statistics}.
Cambridge University Press.\\
\textit{Rigorous treatment of the Bernstein--von Mises theorem (Chapter 10),
establishing that the Bayesian posterior concentrates as
$\mathcal{N}(\hat\theta_\mathrm{MLE},\,\mathcal{I}_T^{-1})$ for large $T$.}

\item[Vehtari et al.\ (2021)]
Vehtari, A., Gelman, A., Simpson, D., Carpenter, B.\ \& B\"urkner, P.-C.\ (2021).
Rank-normalization, folding, and localization: an improved $\hat R$ for assessing
convergence of MCMC.
\textit{Bayesian Analysis}, \textbf{16}(2), 667--718.\\
\textit{Modern definition of MCMC $n_\mathrm{eff}$ and $\hat R$ implemented in
NumPyro's \texttt{summary()} function.}

\end{description}

\end{document}
